\clearpage
\section{Dyadic full characters and the odd critical endpoint}\label{app:dyadic_character}

Source record: Paper\_XI\_Dyadic167. Dependency:
Paper\_XI\_Attachment166.

Evidence: ordinary mathematical proof using the actual Artin--Schreier
matrix tower and the retained Bushnell--Kutzko character and transfer
theorems. No Lean verification is claimed. This note proves no dyadic
quotient-capacity, quadratic-phase, packet-attachment, or H5 result.

\hypertarget{dyadic_character-exact-theorem}{%
\subsection{1. Exact theorem}\label{dyadic_character-exact-theorem}}

Let F=F\_2((t)), let psi\_F be an additive character of conductor t O\_F
(trivial there and nontrivial on O\_F), and let

\begin{Verbatim}[fontsize=\small,breaklines=true,breakanywhere=true]
0<h_1<...<h_d,       each h_i odd,
q_i=3h_i,           N=2^d.
\end{Verbatim}

Use the dyadic Artin--Schreier recursion specified in Section~2. It gives
the actual maximal simple parameter beta\_d with its native arithmetic
defining sequence. There is an actual full simple character theta of its
maximal simple stratum such that

\begin{Verbatim}[fontsize=\small,breaklines=true,breakanywhere=true]
theta restricted to U_F^1 is trivial.                         (1)
\end{Verbatim}

The construction retains the actual canonical beta extension and can be
completed with any prescribed central quasicharacter omega\_F trivial on
U\_F\^{}1. Because \(k_F^\times=\FF_2^\times\) is trivial,
such an omega\_F is unramified. Thus this covers every tamely ramified
central quasicharacter in this particular residue-field setting.

The corresponding compact induction is an actual irreducible admissible
supercuspidal of rank N, by the retained maximal-simple-type theorem. No
claim is made here about its Heisenberg quotient size or about
realization of a supplied dyadic finite packet.

The trace-zero argument from Block 166 does not apply verbatim: at p=2 the
primitive Artin--Schreier parameter has trace one. The correct invariant
is \textbf{integral trace}, not necessarily zero trace. Since psi\_F
kills t O\_F, integral correction trace is sufficient on scalar
principal units. The recursion below preserves the stronger exact value
one.

\hypertarget{dyadic_character-actual-dyadic-recursion-and-exact-native-traces}{%
\subsection{2. Actual dyadic recursion and exact native
traces}\label{dyadic_character-actual-dyadic-recursion-and-exact-native-traces}}

At a recursive step let L be a characteristic-two local field with
residue field F\_2, uniformizer t\_L, and a positive odd q. Define

\begin{Verbatim}[fontsize=\small,breaklines=true,breakanywhere=true]
E=L(alpha),       alpha^2-alpha=t_L^(-q).
\end{Verbatim}

The right side is not an Artin--Schreier coboundary: a negative
valuation of x\^{}2-x is even, whereas -q is odd. The degree-two
polynomial is therefore irreducible and separable. Its valuation
equation gives total ramification and v\_E(alpha)=-q. Put

\begin{Verbatim}[fontsize=\small,breaklines=true,breakanywhere=true]
a=(q+1)/2,       z=t_L^a alpha.
\end{Verbatim}

Then z is an actual uniformizer and its Eisenstein polynomial is

\begin{Verbatim}[fontsize=\small,breaklines=true,breakanywhere=true]
f(Z)=Z^2+t_L^a Z+t_L.                                      (2)
\end{Verbatim}

Its two conjugate alpha values are alpha and alpha+1, so

\begin{Verbatim}[fontsize=\small,breaklines=true,breakanywhere=true]
Tr_(E/L)(alpha)=1.                                         (3)
\end{Verbatim}

Use the coefficient basis 1,z, the trace corestriction pr, and the same
normalized corestriction convention as in Block 166:

\begin{Verbatim}[fontsize=\small,breaklines=true,breakanywhere=true]
kappa=-f'(z)/z = t_L^a/z = alpha^(-1),
s=kappa pr.
\end{Verbatim}

Minus and plus agree in characteristic two. If P0 extracts the constant
coefficient on every E-coordinate, then

\begin{Verbatim}[fontsize=\small,breaklines=true,breakanywhere=true]
u_0=s(P0)=t_L/z^2 is a unit,
Lift_E(C)=u_0^(-1) C P0,
s(Lift_E(C))=C.
\end{Verbatim}

The trace-dual calculation in Block 166 gives, as an identity of L-matrix
traces for every matrix C over E,

\begin{Verbatim}[fontsize=\small,breaklines=true,breakanywhere=true]
tr_L Lift_E(C)
   =Tr_(E/L)(kappa^(-1) tr_E(C))
   =Tr_(E/L)(alpha tr_E(C)).                               (4)
\end{Verbatim}

The lift preserves the complete-chain filtration valuation, by expansion
in 1,z. Thus it is the actual filtered right inverse used in the field
constructor, not an arbitrary inverse on vector spaces.

Define B(L,t\_L;h\_1,\ldots,h\_s) recursively. At s=1 use the L-matrix
of multiplication by alpha with q=q\_1. At s\textgreater=2 use (2) with
q=q\_s, m=2\^{}(s-1), and

\begin{Verbatim}[fontsize=\small,breaklines=true,breakanywhere=true]
C=B(E,z;h_1,...,h_(s-1)),
B(L,t_L;h_1,...,h_s)=alpha I_m+Lift_E(C).                    (5)
\end{Verbatim}

The base trace is one by (3). If tr\_E(C)=1, then (4) gives lift trace
one. Since m is even, the trace of alpha I\_m as an L-matrix is m
Tr\_(E/L)(alpha)=0. Consequently

\begin{Verbatim}[fontsize=\small,breaklines=true,breakanywhere=true]
tr_L B(L,t_L;h_1,...,h_s)=1                                (6)
\end{Verbatim}

for every native recursion, by induction. There is no division by the
rank and no confusion between a native trace and its amplified trace.

The actual field proof from Block 166 applies with this degree-two extension
and normalized lift. At a nontrivial step its hypotheses are
n=2\^{}(s-1)q\_s, R=2\^{}(s-2)q\_(s-1), s(Lift(C))=C, and

\begin{Verbatim}[fontsize=\small,breaklines=true,breakanywhere=true]
n>2R,
-2R>-n-q_1.
\end{Verbatim}

The first strict inequality follows from h\_s\textgreater h\_(s-1).
These are the conditions used in BK Remark 2.2.7. The displayed beta
itself generates the totally ramified degree-2\^{}s field normalizing
the complete chain, and its final critical exponent is -q\_1. This
statement does not identify the final field with its auxiliary derived
field.

For the final length d write the native shorter matrices as

\begin{Verbatim}[fontsize=\small,breaklines=true,breakanywhere=true]
b_j=B(F,t;h_(d-j+1),...,h_d),       1<=j<=d.
\end{Verbatim}

They are the compatible native truncations of (5), because the outer
extension, uniformizer, and lift depend on the same retained largest h
at each stage. For \(1\le j<d\), in rank 2\^{}(j+1), put

\begin{Verbatim}[fontsize=\small,breaklines=true,breakanywhere=true]
gamma_j=b_j tensor I_2,
delta_j=b_(j+1)-gamma_j.
\end{Verbatim}

The defining-sequence calculation of Block 166 gives gamma\_j as a simple
approximation at the critical level r\_j=q\_(d-j) of b\_(j+1), and
delta\_j has grade -r\_j. Equation (6) gives the exact scalar identity

\begin{Verbatim}[fontsize=\small,breaklines=true,breakanywhere=true]
tr_F(gamma_j)=0,
tr_F(delta_j)=1.                                          (7)
\end{Verbatim}

Thus for every scalar u in U\_F\^{}1 on which its additive correction is
defined,

\begin{Verbatim}[fontsize=\small,breaklines=true,breakanywhere=true]
psi_(delta_j)(u I)=psi_F((u-1)tr_F(delta_j))
                 =psi_F(u-1)=1.                           (8)
\end{Verbatim}

It is (8), not a false trace-zero assertion, that drives the full
character construction.

\hypertarget{dyadic_character-the-odd-critical-level-rounding}{%
\subsection{3. The odd-critical-level
rounding}\label{dyadic_character-the-odd-critical-level-rounding}}

The critical levels r=q\_i are odd. Write r=2s+1, and set

\begin{Verbatim}[fontsize=\small,breaklines=true,breakanywhere=true]
m=floor(r/2)=s,        m+1=(r+1)/2.
\end{Verbatim}

The relevant BK statements use the floor, with no assumption that r or
the residue characteristic is odd or even:

\begin{itemize}
\tightlist
\item
  BK 3.1.9(ii) identifies H\^{}k(beta) and H\^{}k(gamma) for
  k\textgreater=floor(r/2)+1.
\item
  BK 3.3.18 gives the correction bijection at
  floor(r/2)\textless=m\textless=r-1.
\item
  BK 3.2.5 extends the resulting character at floor(r/2) to a full
  character, its fibers being the indicated determinant characters.
\end{itemize}

At our exact endpoint m=s, the correction character is legitimate. For
x,y in the radical power P\^{}(s+1),

\begin{Verbatim}[fontsize=\small,breaklines=true,breakanywhere=true]
delta*x*y belongs to P^(2(s+1)-r)=P^1.
\end{Verbatim}

The level-one additive matrix character is trivial on P\^{}1, so the
multiplicative cross term contributes one. Therefore psi\_delta is a
character on U\^{}(s+1), and hence on the recursive subgroup used below.
This checks the borderline odd-level rounding directly. Replacing m by
r/2 as if it were an integer would be incorrect.

The same check applies at the native minimal base level q, which is odd:
its forced additive character begins at floor(q/2)+1=(q+1)/2, and
2*(q+1)/2-q=1.

\hypertarget{dyadic_character-finite-character-extension-at-maximal-rank}{%
\subsection{4. Finite character extension at maximal
rank}\label{dyadic_character-finite-character-extension-at-maximal-rank}}

We reuse the elementary extension procedure with its compatibility
condition explicit. If K/F is a finite totally ramified field, and chi
is a smooth character of U\_K\^{}a trivial on U\_K\^{}a intersect
U\_F\^{}1, define a character on U\_K\^{}a U\_F\^{}1 by

\begin{Verbatim}[fontsize=\small,breaklines=true,breakanywhere=true]
chi_0(ku)=chi(k).
\end{Verbatim}

It is well defined by the intersection condition. Choose M\textgreater=a
with chi\textbar U\_K\^{}M=1. Then chi\_0 is a character of a subgroup
of the finite abelian group \(U_K^1/U_K^M\). Extend it by
adjoining generators and choosing the required complex roots, and pull
it back. The resulting chi\_K extends chi and is trivial on U\_F\^{}1.

All freely chosen character extensions below take place at maximal
native rank, where the centralizer is K itself. We do not apply this
scalar-intersection argument directly inside a nonmaximal matrix
centralizer, which would require a separate determinant-image check.

\hypertarget{dyadic_character-base-full-character-with-trace-one}{%
\subsection{5. Base full character with trace
one}\label{dyadic_character-base-full-character-with-trace-one}}

At native rank two, let E=F{[}alpha{]}, with q=q\_d and a=floor(q/2)+1.
The forced character on U\_E\^{}a is

\begin{Verbatim}[fontsize=\small,breaklines=true,breakanywhere=true]
phi(v)=psi_F(Tr_(E/F)(alpha(v-1))).
\end{Verbatim}

For scalar u in U\_E\^{}a intersect U\_F\^{}1, equation (3) gives

\begin{Verbatim}[fontsize=\small,breaklines=true,breakanywhere=true]
phi(u)=psi_F((u-1)Tr_(E/F)(alpha))=psi_F(u-1)=1.
\end{Verbatim}

Section~4 constructs chi\_E on U\_E\^{}1 extending phi and trivial on
U\_F\^{}1. Define theta\_1 on \(H^1=U_E^1\) U\_(A\_1)\^{}a
by chi\_E on the first factor and psi\_alpha on the second. The
restrictions agree on the intersection. BK 3.2.1 gives the actual full
simple character, with theta\_1\textbar U\_F\^{}1=1.

This uses the true trace value one and only the stated conductor
condition on psi\_F. It does not require an artificial zero-trace
primitive parameter in a quadratic extension.

\hypertarget{dyadic_character-native-rank-induction-via-full-beta-determinant-transfer}{%
\subsection{6. Native-rank induction via full beta-determinant
transfer}\label{dyadic_character-native-rank-induction-via-full-beta-determinant-transfer}}

Assume theta\_j is a full simple character for native b\_j in rank
2\^{}j, with theta\_j\textbar U\_F\^{}1=1. Let K\_j=F{[}b\_j{]}. Since
K\_j is maximal, the character chi\_j=theta\_j\textbar U\_(K\_j)\^{}1 is
its full beta-determinant character.

Transfer theta\_j to the multiplicity-two realization gamma\_j=b\_j
tensor I\_2 in rank 2\^{}(j+1), using BK 3.6.14. Denote the full
transferred character by Theta\_j. BK 3.6.13(ii) preserves the entire
beta-determinant character on U\_(K\_j)\^{}1 at level m0=0. Hence for
every literal scalar u in U\_F\^{}1,

\begin{Verbatim}[fontsize=\small,breaklines=true,breakanywhere=true]
Theta_j(u I)=chi_j(det_(K_j)(u I_2))
            =chi_j(u^2)=theta_j(u^2)=1.                    (9)
\end{Verbatim}

This statement includes nonminimal b\_j. The determinant exponent two is
essential and has not been suppressed.

Now set beta=b\_(j+1), gamma=gamma\_j, delta=delta\_j, r=r\_j,
m=floor(r/2), and

\begin{Verbatim}[fontsize=\small,breaklines=true,breakanywhere=true]
D=H^(m+1)(beta)=H^(m+1)(gamma).
\end{Verbatim}

The equality and endpoint m were verified in Section~3. BK 3.3.18
constructs the actual recursive simple character

\begin{Verbatim}[fontsize=\small,breaklines=true,breakanywhere=true]
vartheta=(Theta_j restricted to D)*psi_delta.
\end{Verbatim}

For u in D intersect U\_F\^{}1, equations (8) and (9) prove

\begin{Verbatim}[fontsize=\small,breaklines=true,breakanywhere=true]
vartheta(u)=1.                                             (10)
\end{Verbatim}

The new field K=F{[}beta{]} is maximal and totally ramified in native
rank 2\^{}(j+1). Its chain has K-period one, so

\begin{Verbatim}[fontsize=\small,breaklines=true,breakanywhere=true]
D intersect U_K^1=U_K^(m+1),
D intersect U_F^1=U_F^ceil((m+1)/2^(j+1)).
\end{Verbatim}

Thus (10) is exactly the compatibility required by Section~4. Extend
vartheta\textbar U\_K\^{}(m+1) to chi\_K on U\_K\^{}1 trivial on
U\_F\^{}1. BK 3.2.5 constructs the full simple character by

\begin{Verbatim}[fontsize=\small,breaklines=true,breakanywhere=true]
theta_(j+1)(kd)=chi_K(k)*vartheta(d),
k in U_K^1, d in D.
\end{Verbatim}

That proposition supplies normalization and well-definedness; its proof
is a character-extension argument and imposes no exclusion of residue
characteristic two. The new theta\_(j+1) is trivial on U\_F\^{}1.
Induction proves (1) in every native rank 2\^{}d.

More generally the character argument needs only integral base field
trace and integral traces for the native corrections. The dyadic tower
establishes their exact value one, while the odd-prime tower of Block 166
establishes zero.

\hypertarget{dyadic_character-canonical-beta-extension-and-prescribed-tame-center}{%
\subsection{7. Canonical beta extension and prescribed tame
center}\label{dyadic_character-canonical-beta-extension-and-prescribed-tame-center}}

Let K/F be the final totally ramified degree-N parameter field and let
eta be the Heisenberg representation for the constructed theta. The
compact type group is

\begin{Verbatim}[fontsize=\small,breaklines=true,breakanywhere=true]
J=O_K^times J^1.
\end{Verbatim}

Here \(\OO_K^\times=U_K^1\) because the residue field is
F\_2. Since U\_K\^{}1 is contained in J\^{}1, we have

\begin{Verbatim}[fontsize=\small,breaklines=true,breakanywhere=true]
J=J^1,       J/J^1=F_2^times={1}.                           (11)
\end{Verbatim}

Thus the actual canonical beta extension on J equals eta; there is no
residue-character choice and no change of that extension. It acts
trivially on the central U\_F\^{}1 by (1). Equation (11) does not
identify any externally prescribed dyadic quadratic phases.

Let omega\_F be any smooth quasicharacter trivial on U\_F\^{}1. Since
\(\OO_F^\times=U_F^1\), it is determined by its value
omega\_F(t). The actual maximal type extends to K\^{}times J.
Concretely, (K\^{}times J)/J is infinite cyclic, generated by a
K-uniformizer. An intertwiner for its conjugation action on eta gives an
extension Lambda\_0. Denote its central character by omega\_0.

Choose a complex scalar s with

\begin{Verbatim}[fontsize=\small,breaklines=true,breakanywhere=true]
s^N=omega_F(t)/omega_0(t).
\end{Verbatim}

Multiply Lambda\_0 by the character of K\^{}times J trivial on J and
taking the K-uniformizer to s. Because t has valuation N in K, this
changes its scalar by s\^{}N and gives the requested central
quasicharacter on all F\^{}times. The value at every element of J,
including the canonical beta extension, is unchanged. This does not
assume that a chosen K-uniformizer raised to N equals t.

The retained maximal-simple-type theorem now gives the actual
irreducible admissible supercuspidal pi=c-Ind\_(K\^{}times
J)\^{}GL\_N(F) Lambda. The noncompact scalar is fixed before any H5
comparison. An incompatible immutably prescribed extension is not
overwritten; as in Block 166, its value must be free or already compatible.

\hypertarget{dyadic_character-what-this-closes-and-what-it-does-not}{%
\subsection{8. What this closes and what it does
not}\label{dyadic_character-what-this-closes-and-what-it-does-not}}

This closes the full-character and actual tame-central-character
construction for the explicit dyadic arithmetic tower. The failure of
trace zero at p=2 is repaired by the correct integral trace invariant.
Both the odd-level floor and the multiplicity-two determinant transfer
are accounted for explicitly.

This note does not compute \(J^1/H^1\) or prove that it
contains a particular dyadic source packet. It does not establish a
quadratic refinement, canonical phase matching, nontrivial stabilizer
action, general transport graph, source admission, or H5. In particular,
an actual supercuspidal host with the required central character is not
asserted to be a completed realization of the original dyadic packet
without those separate proofs.

\hypertarget{dyadic_character-exact-source-checks-and-input-pins}{%
\subsection{9. Exact source checks and input
pins}\label{dyadic_character-exact-source-checks-and-input-pins}}

The following final Block 166 files were read and reused without editing:

\begin{itemize}
\tightlist
\item
  \path{attachment166/central/TRACE_ZERO_CHARACTER_INDUCTION166.md}
  SHA256
  \path{99cb6e80c419e0bf7325e2b2b3fbaba1391903448a3b6a0c4cd9f3944ed406a1}.
\item
  \path{attachment166/tower/TRACE_ZERO_AS_TOWER166.md} SHA256
  \path{49a18fb733f0209b3e38118f5132839a060961edd0bdc0211e1f7b441a4e5653}.
\item
  \path{attachment166/tower/RECURSIVE_TOWER166.md} SHA256
  \path{6f39a077df18eb03f9fab2b6e426529dbe14cf4f7e1cfce95ddb0fcd187d7b33}.
\item
  \path{attachment166/audit/TOWER_CENTRAL_REVIEW166.md} SHA256
  \path{4bc71dd050ddb321a253105c8a4ee7a86185af255532d5605caa7a9de383e26b}.
\end{itemize}

Primary retained source: Bushnell--Kutzko, \emph{The admissible dual of
GL(N) via compact open subgroups}, text SHA256
\path{d69db0827256707a34cd2562af68b9de11fc889ab3f43535ec47b620d58ca9c8}.
The current dyadic check reread Definitions~3.1.7--8 and
Proposition 3.1.9(ii) (printed pp. 90--91), Definition~3.2.1 and
Proposition 3.2.5 (printed pp. 98--100), and Corollary 3.3.18
(printed pp. 112--113). Their level thresholds are floor(r/2)+1 or
floor(r/2), as used above, and their statements do not exclude p=2. The
transfer's preservation of the full beta-determinant character is
BK 3.6.13(ii), with existence Theorem~3.6.14 (printed pp. 136--137), already
directly checked in Block 166 and reread for this application.

The character-extension and correction proofs use neither division by
two nor an odd-order averaging argument. The exact dyadic trace
calculation and the borderline correction-character check are given in
Sections~2--3. No new publication lookup, external
arithmetic-induction/descent constructor, or inverse LLC is used.
